% SOURCE_AUTHORITY: sga5_fr_workpass.tex, lines 8099--8124 (LF-normalized unit).
% SOURCE_SHA256: 791F4EFFC5E02832D5D77ED03518C8156D6F07E4C8238B03545DB93D883FBB28
\subsubsection*{A.3. Uma variante do teorema de Shih}
Suponhamos dado um anel \(S\) filtrado por ideais \(F^i(S)\) \((i\ge0)\), que age sobre \(X\) de forma compatível com sua filtração de sistema projetivo.

Das generalidades de A.1 segue que, para todo inteiro \(i\), \(\gr^i(S)\) atua sobre os objetos bigraduados \(E_r(X)\) por meio de morfismos de bigrau \((i,-i)\). Em particular, para todo inteiro \(n\):
\[
E_r^{(n)}(X)=\bigl(E_r^{p,q}(X)\bigr)_{p+q=n}
\]
é canonicamente dotado de uma estrutura de \((\gr^\bullet(S),\mathcal D)\)-módulo graduado. O mesmo acontece, ao passar ao limite, com
\[
\mathcal E_r^{(n)}(X)=\bigl(\mathcal E_r^{p,q}(X)\bigr)_{p+q=n}.
\]
O anel \(S\), ao agir sobre \(X\) de forma compatível com sua filtração de sistema projetivo, também age sobre \(T^n(X)\) para todo inteiro \(n\). Por conseguinte, o graduado estrito \(\grs(X)\) é canonicamente dotado de uma estrutura de \((\gr^\bullet(S),\mathcal C)\)-módulo graduado e, do mesmo modo, quando \(T^n(X)\) verifica (ML), o graduado estrito \(\grs^\bullet T^n(X)\) é canonicamente dotado de uma estrutura de \((\gr^\bullet(S),\mathcal D)\)-módulo graduado.

\begin{statement}{Teorema A.3}
Suponhamos que, para um inteiro \(n\), os \((\gr^\bullet(S),\mathcal D)\)-módulos graduados
\[
\mathcal E_1^{(n)}(X)=T^n(\grs^\bullet(X))
\quad\text{e}\quad
\mathcal E_1^{(n-1)}(X)=T^{n-1}(\grs^\bullet(X))
\]
são noetherianos. Então:
\begin{enumerate}[label=\textup{(\roman*)}]
\item \(T^n(X)\) verifica (MLAR);
\item \(\grs^\bullet T^n(X)\) é um \((\gr^\bullet(S),\mathcal D)\)-módulo noetheriano.
\end{enumerate}
\end{statement}
